%%% FIELD proof-sharp-attainment
Let \(m=|\Omega|\) and \(k=|\mathcal O|\), enumerate \(\Omega=\{\omega_1,\ldots,\omega_m\}\) and \(\mathcal O=\{o_1,\ldots,o_k\}\), and identify a full law and an observed law with, respectively, \(q=(q_1,\ldots,q_m)\in\mathbb R^m\) and \(p=(p_1,\ldots,p_k)\in\mathbb R^k\).  The observable-data map is the fixed linear map
\[
 (O_\#q)_j=\sum_{i:O(\omega_i)=o_j}q_i,
 \qquad 1\leq j\leq k.
\tag{1}
\]
Thus the observable object is \(p=O_\#q\), whereas the causal object is the complete-law policy functional \(V_\pi(q)\).  We first verify their relationship rather than defining it into existence.

Fix \(q\in\mathcal K(p)\).  The ordinary finite-chain factorization of \(q\), together with \(\mathrm{ass{:}behavior\text{-}factorization}\), gives on every positive-probability full-history cell
\[
 q(A_t=a\mid H_t^F=h)=\mu_t(a\mid O_t(h)).
\tag{2}
\]
Conditional kernels following zero-probability histories may be chosen arbitrarily.  Replace in this chain factorization only the action kernels in (2) by \(\pi_t(\cdot\mid h)\), leaving the initial-state and all post-action kernels fixed, and denote the resulting interventional law by \(q^\pi\).  Assumption \(\mathrm{ass{:}behavior\text{-}positivity}\) states that every factor \(\mu_t(a\mid g)\) is at least \(\epsilon>0\).  Induction over the finite sequence of decision stages therefore shows that every target-policy prefix having positive \(q^\pi\)-mass has positive \(q\)-mass and that, through stage \(t\), its likelihood ratio is
\[
 \prod_{u=1}^t
 \frac{\pi_u(A_u\mid H_u^F)}{\mu_u(A_u\mid H_u^O)}
 =\rho_t^\pi.
\tag{3}
\]
Indeed, the assertion is trivial before the first action; at a decision node, positivity of \(\mu_u\) preserves absolute continuity and multiplies the ratio by \(\pi_u/\mu_u\), while a shared post-action kernel leaves the ratio unchanged.  Marginalizing all variables after \(R_t\) in the two finite chain factorizations and using (3) gives
\[
 \mathbb E_{q^\pi}[R_t]
 =\sum_{\omega\in\Omega}q(\omega)\rho_t^\pi(\omega)R_t(\omega)
 =\ell_t(q).
\tag{4}
\]
Consequently, by \(\mathrm{def{:}stage\text{-}value}\),
\[
 \mathbb E_{q^\pi}\!\left[\sum_{t=1}^H R_t\right]
 =\sum_{t=1}^H\ell_t(q)=V_\pi(q).
\tag{5}
\]
Thus \(\mathrm{def{:}identified\text{-}set}\) takes the image of a causal full-law functional over the laws compatible with the observed distribution; it is not an observed-law regression functional.

We next use finite-dimensional closed-set algebra.  The simplex
\[
 \Delta(\Omega)=\left\{q\in\mathbb R^m:q_i\geq0\ (1\leq i\leq m),\ \sum_{i=1}^m q_i=1\right\}
\tag{6}
\]
is closed and bounded.  Every cylinder marginal \(q_t(h,a,r,m',s')\) and \(q_t^F(h,a)\) is a finite linear combination of coordinates of \(q\).  Hence (1) gives finitely many linear observed-cell equations, the full-history behavior equations in \(\mathrm{ass{:}behavior\text{-}factorization}\) are linear, and every cross-product equation in \(\mathrm{ass{:}weak\text{-}shadow}\) is polynomial of degree at most two.  All index sets are finite.  It follows from \(\mathrm{def{:}compatible\text{-}fiber}\) that \(\mathcal K(p)\) is a closed semialgebraic subset of (6).  Therefore it is compact; it is nonempty by \(\mathrm{ass{:}model\text{-}compatibility}\).

For fixed \(\pi\) and \(\mu\), define
\[
 c_i=\sum_{t=1}^H\rho_t^\pi(\omega_i)R_t(\omega_i),
 \qquad 1\leq i\leq m.
\tag{7}
\]
The reward alphabet is finite, and \(\mathrm{ass{:}behavior\text{-}positivity}\) makes every denominator in (3) at least \(\epsilon\), so every \(c_i\) is a finite real number.  By \(\mathrm{def{:}stage\text{-}value}\),
\[
 V_\pi(q)=\sum_{i=1}^m c_iq_i,
\tag{8}
\]
which is linear and continuous.  Its graph over the fiber,
\[
 \Gamma_p=\left\{(q,v)\in\mathbb R^{m+1}:q\in\mathcal K(p),\ v-\sum_{i=1}^m c_iq_i=0\right\},
\tag{9}
\]
is semialgebraic.  The Tarski--Seidenberg projection theorem \(\mathrm{lem{:}tarski\text{-}seidenberg\text{-}projection}\), applied to the projection of (9) onto its last coordinate, gives
\[
 \{v:\exists q\ ((q,v)\in\Gamma_p)\}=\Theta_I(p)
\tag{10}
\]
as a semialgebraic subset of \(\mathbb R\).  This qualitative projection step permits arbitrary fixed real coefficients; in the real-algebraic case the effective version is also supplied by \(\mathrm{lem{:}effective\text{-}real\text{-}cad}\).  Equation (10) is compact because it is the continuous image under (8) of the compact set \(\mathcal K(p)\).

Finally, a continuous real-valued function on a nonempty compact set attains its minimum and maximum.  Therefore there are \(q^-,q^+\in\mathcal K(p)\) such that
\[
 V_\pi(q^-)=\min_{q\in\mathcal K(p)}V_\pi(q)=L(p),
 \qquad
 V_\pi(q^+)=\max_{q\in\mathcal K(p)}V_\pi(q)=U(p),
\tag{11}
\]
where the final equalities are \(\mathrm{def{:}endpoints}\).  By \(\mathrm{def{:}compatible\text{-}fiber}\), membership of these witnesses means precisely nonnegativity and normalization, reproduction of every observed cell, satisfaction of the known full-history behavior factors, and satisfaction of every weak-shadow cross-product, including boundary cells.  Hence (11) gives attained sharp endpoints of the full compatible-law image, not merely an outer interval.

%%% FIELD proof-compact-face-pasting
The proof uses the elementary minimum-of-a-sum inequality and its equality condition.  By \(\mathrm{thm{:}sharp\text{-}attainment}\), the nonempty fiber \(\mathcal K(p)\) is compact.  Equation (8) in that theorem also shows that each \(\ell_t\) is continuous and linear, so every stagewise minimum and maximum in \(\mathrm{def{:}stagewise\text{-}projection}\) is finite and attained.

For every \(q\in\mathcal K(p)\) and \(t\in\mathcal T\),
\[
 \ell_t^-(p)\leq \ell_t(q)\leq \ell_t^+(p).
\tag{12}
\]
Summing (12) over the finite stage set and using \(V_\pi(q)=\sum_t\ell_t(q)\), \(\mathrm{def{:}rectangular\text{-}relaxation}\), and \(\mathrm{def{:}endpoints}\) gives
\[
 L^{\mathrm{rect}}(p)\leq L(p)\leq U(p)\leq U^{\mathrm{rect}}(p).
\tag{13}
\]

If \(\mathfrak F^-(p)\neq\varnothing\), choose \(q_0\in\mathfrak F^-(p)\).  By \(\mathrm{def{:}extremizing\text{-}faces}\), \(\ell_t(q_0)=\ell_t^-(p)\) for every \(t\), and hence
\[
 V_\pi(q_0)=\sum_t\ell_t^-(p)=L^{\mathrm{rect}}(p).
\tag{14}
\]
Together with (13), this proves \(L(p)=L^{\mathrm{rect}}(p)\).  Conversely, suppose that equality holds and let \(q^-\) attain \(L(p)\), as supplied by \(\mathrm{thm{:}sharp\text{-}attainment}\).  Define
\[
 d_t=\ell_t(q^-)-\ell_t^-(p)\geq0.
\tag{15}
\]
Then
\[
 0=L(p)-L^{\mathrm{rect}}(p)=\sum_{t\in\mathcal T}d_t.
\tag{16}
\]
Because (16) is a finite sum of nonnegative real numbers, every \(d_t=0\).  Thus \(q^-\in F_t^-(p)\) for every \(t\), and \(q^-\in\mathfrak F^-(p)\).  We have proved
\[
 L(p)=L^{\mathrm{rect}}(p)
 \quad\Longleftrightarrow\quad
 \mathfrak F^-(p)\neq\varnothing.
\tag{17}
\]

Apply the same argument to an attained maximizer \(q^+\), using the nonnegative quantities \(d_t^+=\ell_t^+(p)-\ell_t(q^+)\), to obtain
\[
 U(p)=U^{\mathrm{rect}}(p)
 \quad\Longleftrightarrow\quad
 \mathfrak F^+(p)\neq\varnothing.
\tag{18}
\]
By \(\mathrm{def{:}rectangularity\text{-}defect}\), (13) implies \(\Delta^-(p)\geq0\) and \(\Delta^+(p)\geq0\).  If the corresponding intersection in (17) or (18) is empty, the respective endpoint equality is false; the associated weak inequality in (13) is then strict.  This proves \(\Delta^-(p)>0\) or \(\Delta^+(p)>0\), respectively.

%%% FIELD proof-sanity-chen-reduction
The proof uses cross-product algebra to obtain an odds parameterization, an inverse construction to establish sharpness, and importance-ratio cancellation to calculate the target value.  Put
\[
 c_{sa}=p(S_1=s,A_1=a).
\tag{19}
\]
Assumption \(\mathrm{ass{:}model\text{-}compatibility}\) supplies \(\bar q\in\mathcal K(p)\).  Because \(H=1\), \(H_1^F=H_1^O=S_1\).  Observed marginalization and the corrected full-history equation in \(\mathrm{ass{:}behavior\text{-}factorization}\) give
\[
 c_{sa}=\bar q_1^F(s,a)
 =\mu_1(a\mid s)\sum_{a'\in\mathcal A}\bar q_1^F(s,a')
 =\mu_1(a\mid s)p(S_1=s).
\tag{20}
\]
If \(p(S_1=s)>0\), \(\mathrm{ass{:}behavior\text{-}positivity}\) yields
\[
 c_{sa}\geq\epsilon p(S_1=s)>0.
\tag{21}
\]
Thus every action above a positive-mass state belongs to \(\mathcal X_{SA}(p)\), while every cell above a zero-mass state has zero mass under each compatible law.

Fix \((s,a)\in\mathcal X_{SA}(p)\) and \(q\in\mathcal K(p)\).  The division in the following definitions is justified by (21):
\[
 a_{sa}(s',r)=q(M_1=1,S_2=s',R_1=r\mid S_1=s,A_1=a),
\quad
 b_{sa}(s',r)=q(M_1=0,S_2=s',R_1=r\mid S_1=s,A_1=a).
\tag{22}
\]
By \(O_\#q=p\) and \(\mathrm{def{:}static\text{-}chen\text{-}cell\text{-}program}\),
\[
 a_{sa}(s',r)=A_{s,a}[s',r],
 \qquad
 \sum_{r\in\mathcal R}b_{sa}(s',r)=\boldsymbol\beta_{s,a}[s'].
\tag{23}
\]
Dividing the unconditional cross-products in \(\mathrm{ass{:}weak\text{-}shadow}\) by \(c_{sa}^2>0\) gives, for all \(u,s'\in\mathcal S\) and \(r\in\mathcal R\),
\[
 a_{sa}(u,r)b_{sa}(s',r)=b_{sa}(u,r)a_{sa}(s',r).
\tag{24}
\]
Define
\[
 d_{sa,r}=\sum_{u\in\mathcal S}a_{sa}(u,r),
 \qquad
 e_{sa,r}=\sum_{u\in\mathcal S}b_{sa}(u,r).
\tag{25}
\]
The proposition's cellwise response-positivity condition and (23) imply \(d_{sa,r}>0\).  Summing (24) over \(u\) and dividing only by this positive column total gives
\[
 b_{sa}(s',r)=a_{sa}(s',r)w_{sa,r},
 \qquad
 w_{sa,r}=e_{sa,r}/d_{sa,r}\geq0.
\tag{26}
\]
No individual entry of \(A_{s,a}\) is divided by; in particular, (26) forces the matching nonrespondent entry to vanish when a respondent entry is zero.  Summing (26) over \(r\) and using (23) gives
\[
 A_{s,a}w_{sa}=\boldsymbol\beta_{s,a},
\tag{27}
\]
so \(w_{sa}\in\mathcal W_{s,a}\).  Moreover, (26) shows that \(w_{sa}\) uniquely determines, and is uniquely determined by, the conditional completion in this cell.  Hence there is an injective map
\[
 \Phi:\mathcal K(p)\longrightarrow
 \prod_{(s,a)\in\mathcal X_{SA}(p)}\mathcal W_{s,a}.
\tag{28}
\]

For surjectivity, choose any family \((w_{sa})\) in the product in (28), and on each positive cell define
\[
 q^w(R_1=r,M_1=1,S_2=s'\mid S_1=s,A_1=a)=A_{s,a}[s',r],
\tag{29}
\]
\[
 q^w(R_1=r,M_1=0,S_2=s'\mid S_1=s,A_1=a)=A_{s,a}[s',r]w_{sa,r}.
\tag{30}
\]
Feasibility (27) implies that the total conditional mass in (29)--(30) equals
\[
 \sum_{s',r}A_{s,a}[s',r](1+w_{sa,r})
 =p(M_1=1\mid S_1=s,A_1=a)+\mathbf 1^\top\boldsymbol\beta_{s,a}=1.
\tag{31}
\]
Give this conditional distribution baseline mass \(c_{sa}\), and assign zero mass to cells above states with \(p(S_1=s)=0\).  Equations (29), (30), and (27) reproduce every observed respondent and nonrespondent cell, so \(O_\#q^w=p\).  Equation (20), which is a necessary observable implication of nonempty compatibility, verifies the full-history behavior equation for \(q^w\).  For fixed \((s,a,r)\), the response-one and response-zero vectors over \(s'\) are \(A_{s,a}[s',r]\) and \(A_{s,a}[s',r]w_{sa,r}\).  Therefore, for every \(u,s'\),
\[
 A_{s,a}[u,r]A_{s,a}[s',r]w_{sa,r}
 =A_{s,a}[u,r]w_{sa,r}A_{s,a}[s',r].
\tag{32}
\]
Multiplication by \(c_{sa}^2\) verifies the weak-shadow cross-products, including zero entries and zero baseline cells.  Thus \(q^w\in\mathcal K(p)\), and (26), (29), and (30) show that this construction is inverse to \(\Phi\).  Consequently
\[
 \mathcal K(p)\cong
 \prod_{(s,a)\in\mathcal X_{SA}(p)}\mathcal W_{s,a}.
\tag{33}
\]

It remains to calculate the target functional.  By \(\mathrm{def{:}stage\text{-}value}\), positivity makes the denominator below nonzero, and (20), (26), and (29)--(30) imply
\[
\begin{aligned}
 V_\pi(q)
 &=\sum_{s,a,r,m,s'}q(s,a,r,m,s')
   \frac{\pi_1(a\mid s)}{\mu_1(a\mid s)}r\\
 &=\sum_{(s,a)\in\mathcal X_{SA}(p)}
   \frac{c_{sa}\pi_1(a\mid s)}{\mu_1(a\mid s)}
   \sum_{s',r}rA_{s,a}[s',r](1+w_{sa,r})\\
 &=\sum_{(s,a)\in\mathcal X_{SA}(p)}
   p(S_1=s)\pi_1(a\mid s)\theta_{s,a},
\end{aligned}
\tag{34}
\]
where, by \(\mathrm{def{:}static\text{-}chen\text{-}cell\text{-}program}\),
\[
 \theta_{s,a}=\mathbf 1^\top A_{s,a}
 \operatorname{diag}((r)_{r\in\mathcal R})(w_{sa}+\mathbf 1)
 \in\mathcal J_{s,a}.
\tag{35}
\]
Conversely, every independent choice \(\theta_{s,a}\in\mathcal J_{s,a}\) has a witnessing \(w_{sa}\in\mathcal W_{s,a}\), and (29)--(32) paste all such witnesses into one compatible law.  Combining this fact with \(\mathrm{def{:}identified\text{-}set}\) proves the exact image
\[
 \Theta_I(p)=\left\{
 \sum_{(s,a)\in\mathcal X_{SA}(p)}p(S_1=s)\pi_1(a\mid s)\theta_{s,a}:
 \theta_{s,a}\in\mathcal J_{s,a}
 \right\}.
\tag{36}
\]

The extrema in (36) are attained.  Indeed, (27) implies
\[
 \sum_{r\in\mathcal R}d_{sa,r}w_{sa,r}
 =\mathbf 1^\top\boldsymbol\beta_{s,a}.
\tag{37}
\]
Since every \(d_{sa,r}>0\), each feasible coordinate obeys
\[
 0\leq w_{sa,r}\leq
 \frac{\mathbf 1^\top\boldsymbol\beta_{s,a}}{d_{sa,r}}.
\tag{38}
\]
Thus each nonempty \(\mathcal W_{s,a}\) is a compact convex polytope and \(\mathcal J_{s,a}\) is a compact interval.  The weights in (36) are nonnegative, so independent cellwise extremization gives
\[
 L(p)=\sum_{(s,a)\in\mathcal X_{SA}(p)}p(S_1=s)\pi_1(a\mid s)\min\mathcal J_{s,a},
 \qquad
 U(p)=\sum_{(s,a)\in\mathcal X_{SA}(p)}p(S_1=s)\pi_1(a\mid s)\max\mathcal J_{s,a}.
\tag{39}
\]
These are sharp bounds, because every simultaneous choice was realized in (29)--(32).

Finally, if \(\operatorname{rank}(A_{s,a})=|\mathcal R|\) in every positive cell, then necessarily \(|\mathcal S|\geq|\mathcal R|\).  Nonemptiness of \(\mathcal K(p)\) and (28) make every \(\mathcal W_{s,a}\) nonempty.  Proposition \(\mathrm{prop{:}point\text{-}id\text{-}mean}\) therefore makes each \(\mathcal W_{s,a}\), and hence each \(\mathcal J_{s,a}\), a singleton.  Equation (36) is then a singleton.  Under \(X=(S_1,A_1)\), \(F=S_2\), \(Y=R_1\), and response indicator \(M_1\), equations (27), (35), and (39) are the advertised static Chen cell program and its full-column-rank implication.

%%% FIELD proof-fixed-law-value-image
The technique is exact sign-invariant cylindrical algebraic decomposition, using \(\mathrm{lem{:}effective\text{-}real\text{-}cad}\) as the classical algorithmic result and deriving its specialization to this compatible-law fiber.

Fix rational or real-algebraic \(p\), \(\mu\), and \(\pi\).  By \(\mathrm{ass{:}behavior\text{-}positivity}\), every coordinate \(\mu_t(a\mid g)\) appearing in an importance ratio is at least \(\epsilon>0\).  Thus its reciprocal exists and is real algebraic.  The field generated over \(\mathbb Q\) by the finitely many coordinates of \(p\), \(\mu\), and \(\pi\), together with these reciprocals, is an effectively represented real-algebraic coefficient field.

Write \(q=(q_\omega)_{\omega\in\Omega}\).  By \(\mathrm{def{:}compatible\text{-}fiber}\), the formula \(q\in\mathcal K(p)\) is the finite conjunction of
\[
 q_\omega\geq0,
 \qquad \sum_{\omega\in\Omega}q_\omega=1,
 \qquad
 \sum_{\omega:O(\omega)=o}q_\omega=p(o)\quad(o\in\mathcal O),
\tag{40}
\]
the linear full-history behavior equations
\[
 q_t^F(h,a)-\mu_t(a\mid O_t(h))
 \sum_{a'\in\mathcal A}q_t^F(h,a')=0,
\tag{41}
\]
and the quadratic weak-shadow cross-products
\[
 q_t(h,a,r,1,s)q_t(h,a,r,0,s')
 -q_t(h,a,r,0,s)q_t(h,a,r,1,s')=0.
\tag{42}
\]
All cylinder quantities in (41)--(42) are linear sums of the \(q_\omega\).  Moreover, \(\mathrm{def{:}stage\text{-}value}\) gives the linear polynomial equation
\[
 v-\sum_{\omega\in\Omega}q_\omega
 \sum_{t=1}^H\rho_t^\pi(\omega)R_t(\omega)=0
\tag{43}
\]
for \(v=V_\pi(q)\).  Thus
\[
 \Phi(v)\;\Longleftrightarrow\;
 \exists(q_\omega)_{\omega\in\Omega}
 \bigl[q\in\mathcal K(p)\wedge v=V_\pi(q)\bigr]
\tag{44}
\]
is a first-order real formula over that coefficient field.  In the presentation fixed in the theorem, (44) has \(N\) real variables including \(v\), \(s\) polynomial equalities and weak inequalities, and maximum total degree \(d\).

Apply \(\mathrm{lem{:}effective\text{-}real\text{-}cad}\) to all polynomials in (40)--(43), with \(v\) ordered as the uneliminated coordinate.  The algorithm terminates and produces an equivalent quantifier-free univariate sign formula for \(\Phi(v)\), using exact coefficient arithmetic.  The cited CAD bound gives at most
\[
 (sd)^{2^{O(N)}}
\tag{45}
\]
arithmetic operations.  Since producing or inspecting a CAD cell requires at least one algorithmic operation, the number of cells examined is bounded by the same asymptotic quantity.

By \(\mathrm{def{:}identified\text{-}set}\), the truth set of (44) is exactly \(\Theta_I(p)\).  A univariate quantifier-free sign formula changes truth value only at a real root of one of its finitely many nonzero projection polynomials.  Its truth set is therefore a finite union of points and intervals with real-algebraic boundary points.  By \(\mathrm{thm{:}sharp\text{-}attainment}\), \(\Theta_I(p)\) is compact, hence closed and bounded.  After adjacent true CAD cells are merged, every connected component is consequently either a singleton or a closed interval, and all endpoints are real algebraic.

For the exact output, CAD records each boundary root by a defining polynomial and exact root-identification data, such as a Thom encoding.  For the tolerance alternative, let \(E\) be the finite set of distinct component endpoints.  Exact comparison of real-algebraic numbers orders \(E\).  Because \(\eta>0\) and distinct elements of \(E\) have positive separation, repeated rational bisection of their certified isolating intervals yields pairwise disjoint rational intervals, each of width at most \(\eta\).  The defining polynomial and root-identification data certify containment even for an even-multiplicity endpoint, for which a sign-change test would fail.  Pairing the two endpoint enclosures describes each interval component; the single enclosure describes a singleton component.  This proves the stated certified \(\eta\)-description of all components and endpoints.  The procedure eliminates the entire feasible formula (40)--(44), so its certificate is global and does not rely on convexity or on the output of a local optimizer.

%%% FIELD statement-tarski-seidenberg-projection
If \(S\subseteq\mathbb R^{m+k}\) is semialgebraic, then its coordinate projection
\[
\{y\in\mathbb R^k:\exists x\in\mathbb R^m,\ (x,y)\in S\}
\]
is semialgebraic.

%%% FIELD statement-sharp-attainment
For every \(p\) satisfying \(\mathcal K(p)\ne\varnothing\), the fiber \(\mathcal K(p)\) is a compact semialgebraic subset of \(\Delta(\Omega)\), the image \(\Theta_I(p)\) is a compact semialgebraic subset of \(\mathbb R\), and there exist \(q^-,q^+\in\mathcal K(p)\) with \(V_\pi(q^-)=L(p)\) and \(V_\pi(q^+)=U(p)\). Every reported endpoint is therefore attained by a full trajectory law reproducing the observed cells, the known behavior factors, and all weak-shadow cross-product equations.
